\documentclass[10pt]{extarticle}
\usepackage[a4paper,margin=1.5cm]{geometry}
\usepackage[T1]{fontenc}
\usepackage{amsmath,amssymb}
\usepackage{graphicx}
\usepackage[export]{adjustbox}
\usepackage{multicol}
\usepackage{enumitem}
\usepackage{booktabs,array,tabularx}
\setlength{\columnsep}{0.7cm}
\setlength{\parindent}{0pt}
\setlength{\parskip}{3pt}
\pagestyle{empty}
\makeatletter
% figure registry (figdefs.tex is generated by the extract step)
\newcommand{\qfigdef}[4]{\@namedef{qf@file@#1}{#2}\@namedef{qf@w@#1}{#3}\@namedef{qf@wide@#1}{#4}}
\newcommand{\qf@check}[1]{\@ifundefined{qf@file@#1}{\PackageError{score2}{Unknown figure id #1}{}}{}}
\newcommand{\qf@img}[4][]{\includegraphics[width=#2\dimexpr\@nameuse{qf@w@#4}pt\relax,max width=#3,max height=0.45\textheight,#1]{\@nameuse{qf@file@#4}}}
% Full-width questions: finalize wraps every question that contains a wide figure in \qfullbegin...\qfullend,
% so the WHOLE question (stem, figures, options) is set across the page in reading order, then the two columns
% resume. The question's head - its number, stem and first figure (\qfullsplit marks the end of the head) -
% is kept on one page, so a number is never left at a page foot with its figure overleaf; the rest may break
% across pages normally, so no big gaps are left. Without a \qfullsplit, the whole question is kept together
% if it is under 45% of a page. Never write these yourself.
\newif\ifqf@full
\newif\ifqf@boxed
\newsavebox{\qf@box}
\newcommand{\qf@open}{\global\qf@fulltrue\par\global\qf@boxedtrue
  \setbox\qf@box\vbox\bgroup\hsize\textwidth\linewidth\textwidth}
\newcommand{\qf@emit}[1]{\par\egroup\global\qf@boxedfalse
  \ifdim\dimexpr\ht\qf@box+\dp\qf@box\relax<#1\textheight\noindent\usebox{\qf@box}\par\else\unvbox\qf@box\fi}
\newcommand{\qfullsplit}{\ifqf@boxed\qf@emit{0.85}\fi}
\newcommand{\qf@close}{\ifqf@boxed\qf@emit{0.45}\fi\global\qf@fullfalse}
\newcommand{\qfullbegin}{\end{multicols}\qf@open}
\newcommand{\qfullend}{\qf@close\begin{multicols}{2}}
\newcommand{\qfullstart}{\qf@open}
\newcommand{\qfulldone}{\qf@close}
% \qfig[scale]{id}: block figure, centred, directly where it is written (size set by the tool).
% Fallback if a wide figure is not inside a full-width question: break out of the columns in place - no floats.
\newcommand{\qfig}[2][1]{\qf@check{#2}%
  \ifqf@full\qf@blk{#1}{#2}\else
  \ifnum\@nameuse{qf@wide@#2}=1 \expandafter\@firstoftwo\else\expandafter\@secondoftwo\fi
  {\end{multicols}\begin{center}\qf@img{#1}{\textwidth}{#2}\end{center}\begin{multicols}{2}}%
  {\qf@blk{#1}{#2}}\fi}
\newcommand{\qf@blk}[2]{\par\smallskip{\centering\qf@img{#1}{\linewidth}{#2}\par}\smallskip}
% \qfigi[scale]{id}: inline figure (inside option lists / tables / minipages); never breaks columns.
\newcommand{\qfigi}[2][1]{\qf@check{#2}\qf@img[valign=c]{#1}{\linewidth}{#2}}
\makeatother
% structure macros - lint reads these to check completeness and figure placement
\newcommand{\papersec}[1]{\par\bigskip{\centering\Large\bfseries PAPER-#1\par}\medskip}
\newcommand{\partsec}[2]{\par\medskip\noindent{\large\bfseries PART-#1: #2}\par\smallskip}
\newcommand{\sect}[1]{\par\smallskip\noindent\textbf{#1}\par}
\newcommand{\qn}[1]{\par\smallskip\noindent\textbf{#1.}~}
\newcommand{\paraq}[2]{\par\smallskip\noindent\textbf{Paragraph for Questions #1 and #2}\par}
\newenvironment{qoptions}{\begin{itemize}[leftmargin=1.8em,itemsep=1pt,topsep=2pt,label={}]}{\end{itemize}}
\newcommand{\opt}[2]{\item[(#1)] #2}
% figure registry (generated by score2_tools.py finalize)
\qfigdef{p1_phy_s1_q01_1}{TMAS - Test 02 (06 Aug 2023)-img/p1_phy_s1_q01_1.png}{129.4}{0}
\qfigdef{p1_phy_s1_q02_1}{TMAS - Test 02 (06 Aug 2023)-img/p1_phy_s1_q02_1.png}{190.6}{0}
\qfigdef{p1_phy_s1_q03_1}{TMAS - Test 02 (06 Aug 2023)-img/p1_phy_s1_q03_1.png}{158.4}{0}
\qfigdef{p1_phy_s1_q04_1}{TMAS - Test 02 (06 Aug 2023)-img/p1_phy_s1_q04_1.png}{124.2}{0}
\qfigdef{p1_phy_s2_q05_1}{TMAS - Test 02 (06 Aug 2023)-img/p1_phy_s2_q05_1.png}{116.5}{0}
\qfigdef{p1_phy_s2_q06_1}{TMAS - Test 02 (06 Aug 2023)-img/p1_phy_s2_q06_1.png}{106.8}{0}
\qfigdef{p1_phy_s2_q07_1}{TMAS - Test 02 (06 Aug 2023)-img/p1_phy_s2_q07_1.png}{102.1}{0}
\qfigdef{p1_phy_s2_q08_1}{TMAS - Test 02 (06 Aug 2023)-img/p1_phy_s2_q08_1.png}{123.7}{0}
\qfigdef{p1_phy_s2_q09_1}{TMAS - Test 02 (06 Aug 2023)-img/p1_phy_s2_q09_1.png}{145.2}{0}
\qfigdef{p1_phy_s2_q10_1}{TMAS - Test 02 (06 Aug 2023)-img/p1_phy_s2_q10_1.png}{98.0}{0}
\qfigdef{p1_phy_s3_para1_2_1}{TMAS - Test 02 (06 Aug 2023)-img/p1_phy_s3_para1_2_1.png}{97.9}{0}
\qfigdef{p1_phy_s3_para3_4_1}{TMAS - Test 02 (06 Aug 2023)-img/p1_phy_s3_para3_4_1.png}{134.9}{0}
\qfigdef{p1_phy_s3_para5_6_1}{TMAS - Test 02 (06 Aug 2023)-img/p1_phy_s3_para5_6_1.png}{153.9}{0}
\qfigdef{p1_phy_s4_q07_1}{TMAS - Test 02 (06 Aug 2023)-img/p1_phy_s4_q07_1.png}{130.3}{0}
\qfigdef{p1_phy_s4_q08_1}{TMAS - Test 02 (06 Aug 2023)-img/p1_phy_s4_q08_1.png}{136.9}{0}
\qfigdef{p1_phy_s4_q09_1}{TMAS - Test 02 (06 Aug 2023)-img/p1_phy_s4_q09_1.png}{108.2}{0}
\qfigdef{p1_chem_s1_q03_1}{TMAS - Test 02 (06 Aug 2023)-img/p1_chem_s1_q03_1.png}{124.6}{0}
\qfigdef{p1_chem_s1_q03_A}{TMAS - Test 02 (06 Aug 2023)-img/p1_chem_s1_q03_A.png}{28.7}{0}
\qfigdef{p1_chem_s1_q03_B}{TMAS - Test 02 (06 Aug 2023)-img/p1_chem_s1_q03_B.png}{29.7}{0}
\qfigdef{p1_chem_s1_q03_C}{TMAS - Test 02 (06 Aug 2023)-img/p1_chem_s1_q03_C.png}{28.7}{0}
\qfigdef{p1_chem_s1_q03_D}{TMAS - Test 02 (06 Aug 2023)-img/p1_chem_s1_q03_D.png}{28.7}{0}
\qfigdef{p1_chem_s1_q04_1}{TMAS - Test 02 (06 Aug 2023)-img/p1_chem_s1_q04_1.png}{187.5}{0}
\qfigdef{p1_chem_s2_q06_1}{TMAS - Test 02 (06 Aug 2023)-img/p1_chem_s2_q06_1.png}{277.3}{1}
\qfigdef{p1_chem_s2_q06_2}{TMAS - Test 02 (06 Aug 2023)-img/p1_chem_s2_q06_2.png}{27.0}{0}
\qfigdef{p1_chem_s2_q07_1}{TMAS - Test 02 (06 Aug 2023)-img/p1_chem_s2_q07_1.png}{131.1}{0}
\qfigdef{p1_chem_s2_q08_1}{TMAS - Test 02 (06 Aug 2023)-img/p1_chem_s2_q08_1.png}{120.1}{0}
\qfigdef{p1_chem_s2_q08_2}{TMAS - Test 02 (06 Aug 2023)-img/p1_chem_s2_q08_2.png}{44.9}{0}
\qfigdef{p1_chem_s2_q08_3}{TMAS - Test 02 (06 Aug 2023)-img/p1_chem_s2_q08_3.png}{51.4}{0}
\qfigdef{p1_chem_s2_q08_4}{TMAS - Test 02 (06 Aug 2023)-img/p1_chem_s2_q08_4.png}{48.2}{0}
\qfigdef{p1_chem_s2_q08_5}{TMAS - Test 02 (06 Aug 2023)-img/p1_chem_s2_q08_5.png}{28.0}{0}
\qfigdef{p1_chem_s3_q04_1}{TMAS - Test 02 (06 Aug 2023)-img/p1_chem_s3_q04_1.png}{334.6}{1}
\qfigdef{p2_phy_s1_q03_1}{TMAS - Test 02 (06 Aug 2023)-img/p2_phy_s1_q03_1.png}{108.3}{0}
\qfigdef{p2_phy_s1_q05_1}{TMAS - Test 02 (06 Aug 2023)-img/p2_phy_s1_q05_1.png}{45.1}{0}
\qfigdef{p2_phy_s1_q06_1}{TMAS - Test 02 (06 Aug 2023)-img/p2_phy_s1_q06_1.png}{66.3}{0}
\qfigdef{p2_phy_s2_para7_8_1}{TMAS - Test 02 (06 Aug 2023)-img/p2_phy_s2_para7_8_1.png}{108.0}{0}
\qfigdef{p2_phy_s2_para9_10_1}{TMAS - Test 02 (06 Aug 2023)-img/p2_phy_s2_para9_10_1.png}{158.5}{0}
\qfigdef{p2_phy_s3_para1_2_1}{TMAS - Test 02 (06 Aug 2023)-img/p2_phy_s3_para1_2_1.png}{139.4}{0}
\qfigdef{p2_phy_s3_para3_4_1}{TMAS - Test 02 (06 Aug 2023)-img/p2_phy_s3_para3_4_1.png}{68.8}{0}
\qfigdef{p2_phy_s3_para5_6_1}{TMAS - Test 02 (06 Aug 2023)-img/p2_phy_s3_para5_6_1.png}{157.0}{0}
\qfigdef{p2_phy_s4_q07_1}{TMAS - Test 02 (06 Aug 2023)-img/p2_phy_s4_q07_1.png}{90.1}{0}
\qfigdef{p2_phy_s4_q08_1}{TMAS - Test 02 (06 Aug 2023)-img/p2_phy_s4_q08_1.png}{125.9}{0}
\qfigdef{p2_phy_s4_q09_1}{TMAS - Test 02 (06 Aug 2023)-img/p2_phy_s4_q09_1.png}{117.5}{0}
\qfigdef{p2_chem_s1_q05_1}{TMAS - Test 02 (06 Aug 2023)-img/p2_chem_s1_q05_1.png}{172.3}{0}
\qfigdef{p2_chem_s1_q06_1}{TMAS - Test 02 (06 Aug 2023)-img/p2_chem_s1_q06_1.png}{65.6}{0}
\qfigdef{p2_chem_s1_q06_A}{TMAS - Test 02 (06 Aug 2023)-img/p2_chem_s1_q06_A.png}{38.5}{0}
\qfigdef{p2_chem_s1_q06_D}{TMAS - Test 02 (06 Aug 2023)-img/p2_chem_s1_q06_D.png}{45.5}{0}
\qfigdef{p2_chem_s2_para7_8_1}{TMAS - Test 02 (06 Aug 2023)-img/p2_chem_s2_para7_8_1.png}{95.8}{0}
\qfigdef{p2_chem_s3_para1_2_1}{TMAS - Test 02 (06 Aug 2023)-img/p2_chem_s3_para1_2_1.png}{134.7}{0}
\qfigdef{p2_chem_s3_q03_1}{TMAS - Test 02 (06 Aug 2023)-img/p2_chem_s3_q03_1.png}{153.8}{0}
\qfigdef{p2_chem_s3_q03_2}{TMAS - Test 02 (06 Aug 2023)-img/p2_chem_s3_q03_2.png}{186.1}{0}
\qfigdef{p2_chem_s3_q04_1}{TMAS - Test 02 (06 Aug 2023)-img/p2_chem_s3_q04_1.png}{76.8}{0}
\qfigdef{p2_chem_s3_para5_6_1}{TMAS - Test 02 (06 Aug 2023)-img/p2_chem_s3_para5_6_1.png}{118.5}{0}
\raggedcolumns
\makeatletter\@beginparpenalty=10000\makeatother
\newcolumntype{L}[1]{>{\raggedright\arraybackslash}p{#1}}
\begin{document}
\begin{center}{\Large\bfseries TMAS --- Test 02 (06 Aug 2023)}\end{center}
\begin{multicols}{2}
\raggedcolumns
\papersec{1}
\partsec{1}{PHYSICS}
\sect{SECTION-I(i)}
\qn{1} An air capacitor having capacity $C_0$ is filled with dielectric as shown in figure. then the new capacitance of system will be
\qfig{p1_phy_s1_q01_1}
\begin{qoptions}
\opt{A}{$C_0$}
\opt{B}{$2C_0$}
\opt{C}{$3C_0$}
\opt{D}{$4C_0$}
\end{qoptions}
\qn{2} Capacitor 3 in circuit shown below is a variable capacitor. Graph shows the electric potential across capacitor 1 versus $C_3$. The electric potential $V$ across the battery and capacitance $C_1$ is :-
\qfig{p1_phy_s1_q02_1}
\begin{qoptions}
\opt{A}{$2V,\ 8\ \mu F$}
\opt{B}{$10V,\ 6\ \mu F$}
\opt{C}{$10V,\ 8\ \mu F$}
\opt{D}{$10V,\ 12\ \mu F$}
\end{qoptions}
\qn{3} A current $i$ is flowing in a conductor $PQRST$ shaped as shown in the figure. The radius of curved part $QRS$ is $r$ and length of straight portions $PQ$ and $ST$ is very large. The magnetic field at the centre $O$ of the curved part is -
\qfig{p1_phy_s1_q03_1}
\begin{qoptions}
\opt{A}{$\dfrac{\mu_0 i}{4\pi r}\left[\dfrac{3\pi}{2}+1\right]\hat k$}
\opt{B}{$\dfrac{\mu_0 i}{4\pi r}\left[\dfrac{3\pi}{2}-1\right]\hat k$}
\opt{C}{$\dfrac{\mu_0 i}{4\pi r}\left[\dfrac{3\pi}{2}+1\right]\left(-\hat k\right)$}
\opt{D}{$\dfrac{\mu_0 i}{4\pi r}\left[\dfrac{3\pi}{2}-1\right]\left(-\hat k\right)$}
\end{qoptions}
\qn{4} Two identical masses $m$ each are connected by a loose string as shown. Masses $A$ and $B$ are moving with velocity $2V$ and $V$ on a smooth horizontal surface. After sufficiently long time.
\qfig{p1_phy_s1_q04_1}
\begin{qoptions}
\opt{A}{system of masses would have lost a kinetic energy of magnitude $\dfrac{mv^2}{4}$ only if collision between $A$ and $B$ is perfectly inelastic}
\opt{B}{system of masses would not loose any kinetic energy if collision between $A$ and $B$ is elastic}
\opt{C}{system of masses must have lost a kinetic energy $\dfrac{mv^2}{4}$ irrespective of nature of collision between $A$ and $B$.}
\opt{D}{system of masses must have lost all its kinetic energy.}
\end{qoptions}
\sect{SECTION-I(ii)}
\qn{5} In the circuit shown in figure the switch $S$ is closed at $t=0$. A long time after closing the switch
\qfig{p1_phy_s2_q05_1}
\begin{qoptions}
\opt{A}{Voltage drop across the capacitor is $E$}
\opt{B}{Current through the battery is $\dfrac{E}{R_1+R_2+R_3}$}
\opt{C}{Energy stored in the capacitor is {\footnotesize$\dfrac12C\!\left(\dfrac{(R_2+R_3)E}{R_1+R_2+R_3}\right)^2$}}
\opt{D}{Current through the resistance $R_4$ becomes zero}
\end{qoptions}
\qn{6} A point mass $m$, having charge $q$ would be projected from $A$ in a gravity free uniform magnetic field : ($C$ and $D$ are in plane of paper)
\qfig{p1_phy_s2_q06_1}
\begin{qoptions}
\opt{A}{If projection velocity is along $AD$, it can never pass through $D$.}
\opt{B}{If projection velocity is along $AD$, it may pass through $D$ for some specific values of initial velocities (magnitudes)}
\opt{C}{If projection velocity is along $AD$, it may pass through $C$ for certain specific values of projection velocities (magnitudes)}
\opt{D}{If projection velocity is along $AD$, it can never pass through $C$.}
\end{qoptions}
\qn{7} In the circuit shown in figure $C_1=C_2=C$ and capacitor $C_1$ is having initial charge $Q_0$. The switch is closed at $t=0$. Then choose the correct option(s):
\qfig{p1_phy_s2_q07_1}
\begin{qoptions}
\opt{A}{The current at time $t$ in the circuit will be $\left(\dfrac{Q_0}{RC}\right)e^{-2t/RC}$}
\opt{B}{Charge on capacitor $C_1$ at time $t$ will be $\dfrac{Q_0}{2}\left(1+e^{-2t/RC}\right)$}
\opt{C}{Total heat loss in the circuit till the current becomes zero is $\dfrac{Q_0^2}{2C}$}
\opt{D}{Total heat loss in the circuit till the current becomes zero is $\dfrac{Q_0^2}{4C}$}
\end{qoptions}
\qn{8} Consider the system of bodies in figure, their masses are specified in the figure, as well as the length of the string connecting bodies 2 and 3, which are at rest at the same height. Body 1 which is released from rest at a height $h$ makes an inelastic collision with the balance pan. Both strings are considered to be inelastic. The gravitational acceleration $g$ is known (the mass of the balance pan will be neglected)
\qfig{p1_phy_s2_q08_1}
\begin{qoptions}
\opt{A}{Just after sticking to pan both bodies (1) and (2) will have speed given by $\dfrac{m}{M+m}\sqrt{2gh}$}
\opt{B}{Just after sticking to pan both bodies (1) and (2) will have speed given by $\dfrac{m}{2M+m}\sqrt{2gh}$}
\opt{C}{The acceleration of body 3, just after body 1 makes an inelastic collision with the balance pan is $4g\dfrac{m^2}{(M+m)^2}\dfrac{h}{\ell}$}
\opt{D}{The acceleration of body 3, just after body 1 makes an inelastic collision with the balance pan is $2g\dfrac{m^2}{(2M+m)^2}\dfrac{h}{\ell}$}
\end{qoptions}
\qn{9} A massless bucket is initially at rest next to one end of a chain that lies in a straight line on the floor, as shown in Fig. The chain has uniform mass density $\lambda$ (kg/m). You push on the bucket (so that it gathers up the chain) with the force $F(t)$ that gives the bucket and whatever chain is inside, a constant acceleration $a$ at all times. `$t$' is time. Friction is absent
\qfig{p1_phy_s2_q09_1}
\begin{qoptions}
\opt{A}{$F(t)$ at time $t$ is $\dfrac32\lambda a^2t^2$}
\opt{B}{Work done by $F(t)$ upto time $t$ is $\dfrac38\lambda a^3t^4$}
\opt{C}{$F(t)$ at time $t$ is $\dfrac38\lambda a^2t^2$}
\opt{D}{None}
\end{qoptions}
\qn{10} In the figure shown a uniform rod of mass $m$ and length $\ell$ is hinged. The rod is released when the rod makes angle $\theta=60^\circ$ with the vertical.
\qfig{p1_phy_s2_q10_1}
\begin{qoptions}
\opt{A}{The angular acceleration of the rod just after release is $\dfrac{3\sqrt3g}{4\ell}$}
\opt{B}{The normal reaction due to the hinge just after release is $\dfrac{\sqrt{19}mg}{8}$}
\opt{C}{The angular velocity of the rod at the instant it becomes vertical is $\sqrt{\dfrac{3g}{2\ell}}$}
\opt{D}{The normal reaction due to the hinge at the instant the rod becomes vertical is $\dfrac74mg$}
\end{qoptions}
\sect{SECTION-II(i)}
\paraq{1}{2}
A thin, uniform rod with negligible mass and length $0.200$ m is attached to the floor by a frictionless hinge at point $P$ (as shown in figure). A horizontal spring with force constant $k=4.80\ \mathrm{Nm^{-1}}$ connects the other end of the rod to a vertical wall. The rod is in a uniform magnetic field $B=0.340$ T directed into the plane of the figure. There is current $I=6.50$ A in the rod, in the direction shown.
\qfig{p1_phy_s3_para1_2_1}
\qn{1} The torque due to the magnetic force on the rod, about an axis passing though $P$ is $N\times10^{-2}$ Nm then value of $N$ is
\qn{2} When the rod is in equilibrium and makes an angle of $53.0^\circ$ with the floor, the compression/elongation in the spring is $N$ cm. Then $N$ is ?
\paraq{3}{4}
Consider a spool of mass $2$kg and inner and outer radii of $1$m and $2$m respectively acted by forces $F=10$ N as shown in figure. Moment of inertia of the spool about an axis passing through $O$ and perpendicular to the plane of paper is $4\ \mathrm{kgm^2}$.
\qfig{p1_phy_s3_para3_4_1}
\qn{3} The frictional force acting on the body if it performs pure rolling, is found to be $\dfrac{\alpha}{\beta}$, where $\alpha$ and $\beta$ are least integer. Find $(\alpha+\beta)$. :-
\qn{4} If ground is smooth, then the total kinetic energy (in J) after 3 seconds will be :
\paraq{5}{6}
A uniform ring of radius $1$m is performing rolling motion on the horizontal surface as shown in the figure. At $t=0$ acceleration of COM of ring is $3\hat i\ \mathrm{m/s^2}$ while angular acceleration \& angular velocity are $-3\hat k\ \mathrm{rad/s^2}$ and $-2\hat k\ \mathrm{rad/s}$ where $x$-$z$ plane is horizontal and $y$ is vertical axis. $A$ is in contact with surface and $C$ is at $2$m above the surface. $B$ and $D$ are $1$m above the surface.
\qfig{p1_phy_s3_para5_6_1}
\qn{5} Magnitude of acceleration (in $\mathrm{m/s^2}$) of point $B$ with respect to $D$ is :
\qn{6} Magnitude of acceleration (in $\mathrm{m/s^2}$) of point $A$ is :
\sect{SECTION-II(ii)}
\qn{7} Sand is transported to the construction site using a conveyor belt. The length of the belt is $\ell$, the angle with respect to the horizontal is $\alpha$ ; the belt is driven by the lower pulley with radius $R$, powered externally. The sand is put onto the belt at a constant rate $\mu$ (kg/s). What is the minimal required torque (in N-m) needed to transport the sand ? The coefficient of friction is larger enough for the sand grains to stop moving immediately after hitting the belt ; take the initial velocity of the sand grains to be zero. (Take $\mu=0.2$ kg/sec, $R=1$m, $\ell=20$ m and $\alpha=30^\circ$)
\qfig{p1_phy_s4_q07_1}
\qn{8} A non conducting rod of length $2$m is hinged at one end and a charge of $\dfrac1{10}$C is distributed uniformly over it. At $t=0$, it is released from the position shown. There exist a uniform magnetic field of $\sqrt{15}$ T inside the plane of motion of rod. If mass of rod is $100$ g, then find the value of hinge force (in N) when rod is rotated by $\dfrac\pi2$ due to gravity.
\qfig{p1_phy_s4_q08_1}
\qn{9} In the figure net magnetic force on wire $AB$ due to infinite fixed wire is given as $F=\dfrac{N\mu_0Ii}{4\pi}\ell n2$, then $N$ is
\qfig{p1_phy_s4_q09_1}
\partsec{2}{CHEMISTRY}
\sect{SECTION-I(i)}
\qn{1} A protein molecule crystallizes as orthorhombic with $a=80$ \AA, $b=30$ \AA{} and $c=125$ \AA. If there are six protein molecules per unit cell and density of unit cell is $1.5$ g/ml then the molar mass of protein in kg/mole is :\\
(Given : $N_A=6\times10^{23}$)
\begin{qoptions}
\opt{A}{45000}
\opt{B}{45}
\opt{C}{15}
\opt{D}{75}
\end{qoptions}
\qn{2} Enthalpy for the reaction :\\
$\mathrm{Ag^+(aq)+Br^-(aq)\to AgBr(s)}$ is $-100$ KJ.\\
If magnitude of enthalpy of formation of $\mathrm{Ag^+(aq)}$ and $\mathrm{Br^-(aq)}$ is in the ratio of $5:6$ while formation of $\mathrm{Ag^+(aq)}$ is an endothermic where as formation of $\mathrm{Br^-}$ is an exothermic process. If enthalpy of formation of $\mathrm{AgBr(s)}=-110\ \mathrm{KJ\,mol^{-1}}$ then the enthalpy of formation of $\mathrm{Ag^+(aq)}$ is :
\begin{qoptions}
\opt{A}{$50\ \mathrm{KJ\,mol^{-1}}$}
\opt{B}{$60\ \mathrm{KJ\,mol^{-1}}$}
\opt{C}{$1000\ \mathrm{KJ\,mol^{-1}}$}
\opt{D}{$10\ \mathrm{KJ\,mol^{-1}}$}
\end{qoptions}
\qn{3} $P\xrightarrow[\mathrm{(iii)\ Excess\ PBr_3,\ (iv)\ Excess\ KCN,\ (v)\ H^+,H_2O}]{\mathrm{(i)\ O_3,\ then\ Zn\ in\ H_2O,\ (ii)\ NaBH_4}}$
\qfig{p1_chem_s1_q03_1}
Product $P$ is
\begin{qoptions}
\opt{A}{\qfigi{p1_chem_s1_q03_A}}
\opt{B}{\qfigi{p1_chem_s1_q03_B}}
\opt{C}{\qfigi{p1_chem_s1_q03_C}}
\opt{D}{\qfigi{p1_chem_s1_q03_D}}
\end{qoptions}
\qn{4} 
\qfig{p1_chem_s1_q04_1}
Incorrect statement is :-
\begin{qoptions}
\opt{A}{`A' can be moulded in variety of shapes.}
\opt{B}{`B' is cast iron}
\opt{C}{`C' contains more impurities than `A' and `B'}
\opt{D}{`B' has about $3\%$ carbon content.}
\end{qoptions}
\sect{SECTION-I(ii)}
\qn{5} For the following mechanism one of the involved species is a catalyst :\\
\textbf{Step (i) :} {\small$\mathrm{H^++HNO_2}\mathrel{\overset{K_1}{\underset{K_2}{\rightleftharpoons}}}\mathrm{H_2NO_2^+}$}\\
\textbf{Step (ii) :} {\small$\mathrm{H_2NO_2^++Br^-}\mathrel{\overset{K_3}{\rightleftharpoons}}\mathrm{NOBr+H_2O}$}\\
\textbf{Step (iii) :}\\ {\small$\mathrm{NOBr+C_6H_5NH_2}\mathrel{\overset{K_4}{\rightleftharpoons}}\mathrm{C_6H_5N_2^++Br^-+H_2O}$}\\
Select the correct option(s) :
\begin{qoptions}
\opt{A}{Intermediates are $\mathrm{H_2NO_2^+}$ and $\mathrm{NOBr}$}
\opt{B}{Catalyst is $\mathrm{Br^-}$}
\opt{C}{If step (ii) forward reaction is slowest step then overall order of reaction is $3$.}
\opt{D}{If step (i) forward reaction is slowest step then overall order of reaction is $2$.}
\end{qoptions}
\qfullbegin
\qn{6} 
\qfig{p1_chem_s2_q06_1}
\qfullsplit
If $G$ gives $\mathrm{CO_2}$ with $\mathrm{NaHCO_3}$ then choose the incorrect statement(s) :
\begin{qoptions}
\opt{A}{Product $C$ is \qfigi{p1_chem_s2_q06_2}}
\opt{B}{$1$ --COOH group is present in product $F$}
\opt{C}{Compound $G$ reacts with alcohol to give fruity smell}
\opt{D}{$F$ gives positive Tollen's test.}
\end{qoptions}
\qfullend
\qn{7} $\mathrm{Ph-CHO}\xrightarrow[\mathrm{(ii)\,H^+/H_2O}]{\mathrm{(i)\,CH_3MgBr}}\mathrm{(A)}\xrightarrow{\mathrm{PCC}}\mathrm{(B)}$ \qfigi{p1_chem_s2_q07_1}\par
Select the correct statement
\begin{qoptions}
\opt{A}{Acidic strength of $E$ is more than $D$}
\opt{B}{Both $D$ and $E$ gives immediate turbidity with Lucas reagent.}
\opt{C}{DOU of $D$ is $4$}
\opt{D}{Product $E$ is ethanol}
\end{qoptions}
\qn{8} Predict product $P$, $Q$, $R$, $S$?
\qfig{p1_chem_s2_q08_1}
\begin{qoptions}
\opt{A}{Product $P$ is \qfigi{p1_chem_s2_q08_2}}
\opt{B}{Product $Q$ is \qfigi{p1_chem_s2_q08_3}}
\opt{C}{Product $R$ is \qfigi{p1_chem_s2_q08_4}}
\opt{D}{Product $S$ is \qfigi{p1_chem_s2_q08_5}}
\end{qoptions}
\qn{9} The correct statement(s) related to the metal extraction processes is(are)
\begin{qoptions}
\opt{A}{A mixture of $\mathrm{PbS}$ and $\mathrm{PbO}$ undergoes self-reduction to produce $\mathrm{Pb}$ and $\mathrm{SO_2}$.}
\opt{B}{In the extraction process of copper from copper pyrites, silica is added to produce copper silicate.}
\opt{C}{Partial oxidation of sulphide ore of copper by roasting, followed by self-reduction produces blister copper.}
\opt{D}{In cyanide process, zinc powder is utilized to precipitate gold from $\mathrm{Na[Au(CN)_2]}$}
\end{qoptions}
\qn{10} Select correct option for given ions.\\
(i) Glycinate ($\mathrm{gly^{-1}}$)\\
(ii) Acetylacetonato ($\mathrm{acac^-}$)\\
(iii) Dipyridine (dipy)\\
(iv) Butylenediamine (bn)\\
(v) Ethylenediamine (en)\\
(vi) O-phenanthroline (O-phen)\\
(vii) Acetate ($\mathrm{Ac^-}$)\\
(viii) $\mathrm{EDTA^{4-}}$ (Ethylenediaminetetraacetate)\\
(ix) Diethylene triamine (dien)
\begin{qoptions}
\opt{A}{Six ligands form 5 membered ring containing atleast 2 nitrogen atoms and metal centre}
\opt{B}{Only one ligand form 5 membered ring containing two oxygen atom and metal centre}
\opt{C}{Seven ligands form 5 membered ring containing 2 carbon atoms and metal centre}
\opt{D}{Three ligands already contain a 6 membered ring.}
\end{qoptions}
\sect{SECTION-II(i)}
\paraq{1}{2}
An ideal binary solution of $A$ and $B$ has a vapour pressure of $900$ torr. It is distilled at constant temperature till $\dfrac23$ of the original amount is distilled. If the mole fraction of `$A$' in the residue and mole fraction of `$B$' in the condensate are $0.3$ and $0.4$ respectively and vapour pressure of residue is $860$ torr then :
\qn{1} The mole fraction of `$A$' in the initial mixture is :
\qn{2} Vapour pressure of pure `$A$' in torr is equal to :
\paraq{3}{4}
\par\noindent\hfill{\footnotesize$\text{Cyclopropyl benzene}\xrightarrow[\Delta]{\mathrm{H^+/KMnO_4}}\mathrm{(A)}\xrightarrow{\mathrm{PCl_5}}\mathrm{(B)}\xrightarrow[\mathrm{NH_4HCl}]{\mathrm{EtMgBr\ (1\ eq.)}}\mathrm{(C)}\xrightarrow{\mathrm{SeO_2}}\mathrm{(D)}\xrightarrow{\mathrm{HIO_4}}P+Q$}\hfill\null\par
In product $Q$, $4\pi$ electrons are involved in resonance
\qn{3} D.O.U. of $P$ is?
\qfullbegin
\qn{4} How many of the following are more acidic than $Q$
\qfig{p1_chem_s3_q04_1}
\qfullsplit
$\mathrm{CH_3CH_2OH}$ , $\mathrm{HCl}$, $\mathrm{CF_3COOH}$
\qfullend
\paraq{5}{6}
Where two or more than two complexes having similar molecular weight (or) molecular formula but their structural arrangements are different these forms are known as isomer \& this phenomenon is known as isomerism; isomerism are of two types structural \& stereoisomerism ; stereoisomerism are of two types geometrical isomerism \& optical isomerism; compound which has plane of symmetry will not exhibit optical isomerism.
\qn{5} Find total number of bidentate ligands ($X$) for which complex $\mathrm{M}(X)_3$ can show geometrical as well as optical isomerism\\
\textbf{en} (ethylenediamine), $\textbf{OX}^{-2}$ (oxalate),\\
\textbf{dipy}(dipyridine), $\textbf{gly}^{-1}$(glycinate),\\
\textbf{pn}(propylenediamine)
\qn{6} Total number of stereoisomers possible for complex $[\mathrm{M(AB)bcde}]$.
\sect{SECTION-II(ii)}
\qn{7} The octahedral complex $[\mathrm{Rh(NO_2)(SCN)(en)_2}]^+$ can exist in a total number of isomeric forms including stereoisomers is $Y$, find $\dfrac Y2$ :-
\qn{8} For the $[\mathrm{Cr(H_2O)_6}]^{2+}$ ion, the mean pairing energy $P$ is found to be $23500\ \mathrm{cm^{-1}}$. The magnitude of $\Delta_0$ is $13900\ \mathrm{cm^{-1}}$. Calculate the C.F.S.E. ($\mathrm{cm^{-1}}$) for this complex ion corresponding to high spin state ($x$) and low spin state ($y$). Write your answer as $\left(\dfrac{y-x}{100}\right)$
\qn{9} An element $X$ (Atomic weight $=24$ gm/mole) form FCC lattice. If edge length of unit cell is $4\times10^{-8}$ cm and observed density is $2.4\times10^3\ \mathrm{kg/m^3}$, then percentage occupancy of lattice points of element $X$ is\\
(Given $N_A=6\times10^{23}$)
\partsec{3}{MATHEMATICS}
\sect{SECTION-I(i)}
\qn{1}\par\noindent\hfill{\small$\displaystyle\int\Bigl(2e^{\tan^{-1}x+x^2+\ln x}$}\hfill\null\par\noindent\hfill{\small$\displaystyle+\dfrac{2x^3+2x+1}{x+\frac1x}e^{\tan^{-1}x+x^2+\ln x}\Bigr)dx$}\hfill\null\par is equal to (where $C$ is constant of integration)
\begin{qoptions}
\opt{A}{$\dfrac{x^3}{3}e^{\tan^{-1}x+x^2+\ln x}+C$}
\opt{B}{$\dfrac{x^2}{2}e^{\tan^{-1}x+x^2+\ln x}+C$}
\opt{C}{$xe^{\tan^{-1}x+x^3+\ln x}+C$}
\opt{D}{$x^2e^{\tan^{-1}x+x^2}+C$}
\end{qoptions}
\qn{2} The function $f(x)=\alpha x^2-\beta x-4x^3+\gamma$, where $\alpha,\beta,\gamma\in R$, has local minima at $P(\log_2a,\ f(\log_2a))$ \& local maxima at $Q(\log_2a^2,\ f(\log_2a^2))$. If the graph of $f(x)$ changes concavity about the point $R\left(\dfrac34,\ f\left(\dfrac34\right)\right)$, then which of the following conic section can have eccentricity `$a$' -
\begin{qoptions}
\opt{A}{circle}
\opt{B}{parabola}
\opt{C}{ellipse}
\opt{D}{hyperbola}
\end{qoptions}
\qn{3} If $m$ is the slope of a line which is tangent to $y^3=x^4$ \& a normal to $x^2-2x+y^2=0$, then $\left(\dfrac{3m}{4}\right)^3$ is equal to ($m\neq0$)
\begin{qoptions}
\opt{A}{$3$}
\opt{B}{$\dfrac43$}
\opt{C}{$4$}
\opt{D}{$\dfrac34$}
\end{qoptions}
\qn{4} For the equation $b\sqrt{1-a^2}-a\sqrt{1-b^2}=\sin\left(\dfrac{b-a}{\sqrt{1-x^2}}\right)$, where $0<a<b<1/2$, then which of the following is true -
\begin{qoptions}
\opt{A}{there is no real root of the equation in $(a,b)$}
\opt{B}{there is exactly one real root of the equation in $(a,b)$}
\opt{C}{there is atleast one root of the equation in $(a,b)$}
\opt{D}{none of these}
\end{qoptions}
\sect{SECTION-I(ii)}
\qn{5} Let $f(x,y)=\dfrac{1}{\displaystyle\lim_{\alpha\to0}\left(\dfrac{\sqrt[x]{\sec\alpha}-\sqrt[y]{\sec\alpha}}{\alpha^2}\right)}$, then
\begin{qoptions}
\opt{A}{$f(2,3)=12$}
\opt{B}{$f(x,y)=12$ has infinitely many solutions}
\opt{C}{Slope of tangent to the curve $f(x,y)=12$ at $x=2$ is $4$}
\opt{D}{$f(x,y)=12$ has no horizontal tangent}
\end{qoptions}
\qn{6} If $\displaystyle\int\cot\left(2\tan^{-1}\sqrt{\dfrac{\sqrt{1+\sqrt x}-x^{1/4}}{\sqrt{1+\sqrt x}+x^{1/4}}}\right)dx=\dfrac{qx^{p/q}}{p}+C$; (where $p$ and $q$ are relatively prime and $C$ is constant of integration), then
\begin{qoptions}
\opt{A}{$p=5$}
\opt{B}{$p+2q=13$}
\opt{C}{$p-q=1$}
\opt{D}{$q=5$}
\end{qoptions}
\qn{7} If $y^x-x^y=1$, then $\dfrac{dy}{dx}$ at $x=1$, is -
\begin{qoptions}
\opt{A}{$2(1-\ell n2)$}
\opt{B}{$2(1+\ell n2)$}
\opt{C}{$\ell n(4e^2)$}
\opt{D}{$\ell n\left(\dfrac{e^2}{4}\right)$}
\end{qoptions}
\qn{8} Identify the CORRECT statements -
\begin{qoptions}
\opt{A}{Sum of two periodic functions may or maynot be a periodic function}
\opt{B}{A differentiable function in $R$ will always be continuous in $R$}
\opt{C}{A continuous function whose domain is in a closed interval is always bounded}
\opt{D}{In $[a,b]$ for a function $f(x)$, if $f(a).f(b)<0$, then there exists atleast one root of $f(x)=0$ in $(a,b)$}
\end{qoptions}
\qn{9} If $\displaystyle\int\dfrac{x+\sqrt[3]{x^2}+\sqrt[6]{x}}{x\left(1+\sqrt[3]{x}\right)}dx=A.x^B+C\tan^{-1}\left(x^D\right)+E$, then which of the following is/are true (where `$E$' is constant of integration) ?
\begin{qoptions}
\opt{A}{$A=\dfrac23$}
\opt{B}{$AB=1$}
\opt{C}{$(AB+C)$ is a prime number}
\opt{D}{Sum of all coefficients of even power of $x$ in the expansion of $\left(AB+\dfrac xD\right)^C$ is equal to $\dfrac{(C+1)^C+(C-1)^C}{2}$}
\end{qoptions}
\qn{10} If $n\in N$ and $f(x)=({}^nC_0x+1)({}^nC_1x+2)({}^nC_2x+3)({}^nC_3x+4)\ldots({}^nC_nx+n+1)$ and $f'(0)=\dfrac32(7!)$, then $n$ is equal to
\begin{qoptions}
\opt{A}{$3$}
\opt{B}{$4$}
\opt{C}{$5$}
\opt{D}{$6$}
\end{qoptions}
\sect{SECTION-II(i)}
\paraq{1}{2}
A Rod of length $10$cm slides between two lines, $x$-axis \& $y=-\sqrt3x$. If locus of mid point of rod is curve `$S$'.
\qn{1} $x$-intercept of $S$ is
\qn{2} If maximum distance between origin and curve $S$ is equal to $k$ then value of $k\sqrt3$ is
\paraq{3}{4}
Let $f:[0,\infty)\to[2,\infty)$ be a differentiable increasing and onto function satisfying\\
$f^2(x)-2f(x)-\sqrt x=f^2(y)-2f(y)-\sqrt y,\ \forall\ x,y\in[0,\infty)$
\qn{3} The value of $\displaystyle\lim_{x\to0^+}\dfrac{f(x)-2}{\sqrt x}$ is equal to -
\qn{4} If $g(x)$ be the inverse of $f(x)$, then $g'(4)$ is equal to -
\paraq{5}{6}
Let $f:R\to R$ be a function defined as $f(x)=3x^5-25x^3+60x+5$ and\\
$g(x)=\begin{cases}\max\{f(t);-4\le t\le x\} & ;\ -4\le x\le0\\ \min\{f(t);0<t\le x\} & ;\ 0<x\le2\\ f(x)-16 & ;\ x>2\end{cases}$
\qn{5} Total number of points at which the graph of $y=f(x)$ and $y=g(x)$ is non-differentiable in $[-4,\infty)$ is/are
\qn{6} If the equation $f(x)+\lambda=0$ is having exactly three distinct real roots then total number of possible integral values of $\lambda$ are
\sect{SECTION-II(ii)}
\qn{7} $\displaystyle\lim_{x\to\infty}x^3\left(\sqrt{x^2+\sqrt{1+x^4}}-\sqrt2x\right)=\dfrac{\sqrt2}{\lambda}$ find value of $\lambda$
\qn{8} Let $P(x)$ be a polynomial of degree $4$ with $P(1)=1$, $P'(1)=0$, $P''(1)=2$, $P'''(1)=-12$ and $P''''(1)=24$ then value of $\dfrac{P(3)}{P(2)}$ is equal to
\qn{9} Let $\displaystyle\int\dfrac{(x^2-1)\,dx}{x^3\sqrt{3x^4+2x^2-1}}=f(x)+c$, (where $c$ is constant of integration) and $f(1)=-1$. If $\lambda=\displaystyle\lim_{x\to\infty}f(x)$, then the absolute value of $\dfrac{2\lambda}{\sqrt3}$ is
\raggedcolumns
\papersec{2}
\partsec{1}{PHYSICS}
\sect{SECTION-I(i)}
\qn{1} A proton is fired from origin with velocity $\vec V=V_0\hat j+V_0\hat k$ in a uniform magnetic field $\vec B=B_0\hat j$. In the subsequent motion of the proton ($v_0$ and $B_0$ are positive)
\begin{qoptions}
\opt{A}{its $z$-coordinate can never be negative}
\opt{B}{its $x$-coordinate can never be positive}
\opt{C}{its $y$-coordinate will be proportional to time}
\opt{D}{None}
\end{qoptions}
\qn{2} A charged particle of specific charge $\alpha$ moves with velocity $\vec v=v_0\hat i$ in a magnetic field $\vec B=\dfrac{B_0}{\sqrt2}\left(\hat j+\hat k\right)$. Then (specific charge = charge per unit mass)
\begin{qoptions}
\opt{A}{Path of the particle is a helix}
\opt{B}{Path of the particle is circle}
\opt{C}{Distance moved by the particle in time $t=\dfrac{\pi}{\alpha B_0}$ is $\dfrac{\pi v_0}{B_0\alpha}$}
\opt{D}{None}
\end{qoptions}
\qn{3} For the given circuit in steady state, choose the correct statements :- (given all capacitors have capacitance = $1$ in SI \& all resistors have resistance = $10$ in SI unit)
\qfig{p2_phy_s1_q03_1}
\begin{qoptions}
\opt{A}{Potential difference between point $A$ \& $B$ $(V_A-V_B)$ is $40$ V.}
\opt{B}{Potential difference between point $A$ \& $B$ $(V_A-V_B)$ is $50$ V.}
\opt{C}{Magnitude of charge $Q_1$ is $30$ C.}
\opt{D}{Magnitude of charge $Q_2$ is $50$ C.}
\end{qoptions}
\qn{4} As preparation for the olympiads, Ivan wanted to make a graphite-filled, parallel plate, $8.85$ nF capacitor with square plates of side $1$ m. ``Such big metal plates will be expensive,'' said his father. ``Can you manage with cardboard? Cover it with aluminium foil to have metal.'' His sister who studies physics in college, said, ``To get the correct capacitance'', make sure you keep the plates parallel. Having parallel plates with a fractional error of $0.001$ in plate separation is better than having the correct plate separation at one end and an angle of $0.001$ radians between the plates.'' The dielectric constant of graphite is $10$ and its dielectric strength is $1$ kV/mm. Which of the following statements are correct?
\begin{qoptions}
\opt{A}{The aluminium foil on cardboard will not work.}
\opt{B}{The plate separation required is $1$ cm.}
\opt{C}{The capacitor can withstand a maximum potential difference of $1000$ V.}
\opt{D}{Ivan's sister is right.}
\end{qoptions}
\qn{5} For technology project, a student has built a vehicle, of total mass $5$ kg, that moves itself. As shown in the figure, it runs on four light wheels. A reel is attached to one of the axles, and a cord originally wound on the reel goes up over the pulley attached to the vehicle to support an elevated load. After the vehicle is released from rest, the load descends very slowly, unwinding the cord to turn the axle and make the vehicle move forward. Friction is negligible in pulley and axle bearings. The wheels do not slip on the floor. The reel has been constructed with a conical shape so that load descends at a constant low speed while the vehicle moves horizontally across the floor with a constant acceleration, reaching a final velocity of $3$ m/s. Mark the correct option(s) :
\qfig{p2_phy_s1_q05_1}
\begin{qoptions}
\opt{A}{Magnitude of impulse on the vehicle due to friction of floor is $15$ N-s.}
\opt{B}{Work done by friction force between tyre and floor is $27$ J.}
\opt{C}{Change in gravitational potential energy of earth and cart system is $-27$ J.}
\opt{D}{The acceleration of the cart is due to friction force from the ground.}
\end{qoptions}
\qn{6} A putty ball of mass $m$ and speed $v_0$ strikes a dumbell as shown in figure. Dumbbell consists of light rod and two particles. Putty sticks to particle after collision. Mark the CORRECT option (s) :-
\qfig{p2_phy_s1_q06_1}
\begin{qoptions}
\opt{A}{Angular velocity of dumbell is $\dfrac{v_0}{4\ell}$ clockwise.}
\opt{B}{Angular velocity of dumbell is $\dfrac{3v_0}{\ell}$ clockwise.}
\opt{C}{Energy lost in collision is $\dfrac{5mv_0^2}{16}$}
\opt{D}{Energy lost in collision is $\dfrac{3mv_0^2}{16}$}
\end{qoptions}
\sect{SECTION-I(ii)}
\paraq{7}{8}
A system consists of two small spheres of mass $1$kg and $2$kg interconnected by a weightless spring. At $t=0$, the spheres are set in motion with initial velocities as shown in figure. Position vector of centre of mass at $t=0$ is $20\,\hat j$. Initially spring is in its natural length. System moves in earth's gravitational field.
\qfig{p2_phy_s2_para7_8_1}
\qn{7} The momentum of the system at $t=2$ sec in kg m/s is $\left(x\hat i+y\hat j\right)$. Find $(x+y)$ ($g=10\ \mathrm{m/s^2}$)
\begin{qoptions}
\opt{A}{$-46$}
\opt{B}{$-92$}
\opt{C}{$80$}
\opt{D}{None of these}
\end{qoptions}
\qn{8} The position vector of centre of mass of the system at $t=1$ s is $\left(x\hat i+y\hat j\right)$. Find $(x+y)$.
\begin{qoptions}
\opt{A}{$19.67$}
\opt{B}{$40.3$}
\opt{C}{$67.8$}
\opt{D}{None}
\end{qoptions}
\paraq{9}{10}
At $t=0$, switch is closed. Capacitors are initially uncharged. Current through battery at $t=0^+$ is $(10\alpha)$ Ampere while that at $t\to\infty$ is $(10\beta)$ Ampere.
\qfig{p2_phy_s2_para9_10_1}
\qn{9} The value of $\alpha$ is \underline{\hspace{1.5cm}}
\begin{qoptions}
\opt{A}{$2$}
\opt{B}{$1.5$}
\opt{C}{$4.5$}
\opt{D}{$1$}
\end{qoptions}
\qn{10} Value of $\beta$ is \underline{\hspace{1.2cm}}
\begin{qoptions}
\opt{A}{$0.3$}
\opt{B}{$3$}
\opt{C}{$1$}
\opt{D}{$2$}
\end{qoptions}
\sect{SECTION-II(i)}
\paraq{1}{2}
Two charge particles each of mass `$m$', carrying charge $+q$ and connected with each other by a massless inextensible string of length $2L$ are describing circular path in the plane of paper, each with speed $v=\dfrac{qB_0L}{m}$ (where $B_0$ is constant) about their centre of mass in the region in which an uniform magnetic field $\vec B$ exists into the plane of paper as shown in figure. Neglect any effect of electrical \& gravitational forces.
\qfig{p2_phy_s3_para1_2_1}
\qn{1} The magnitude of the magnetic field such that no tension is developed in the string will be $\dfrac{B_0}{N}$, find $N$ :-
\qn{2} If the actual magnitude of magnetic field is half to that of calculated in part (i) then tension in the string is $\dfrac{q^2B_0^2L}{Nm}$, where $N$ is least integer. Find the value of $N$.
\paraq{3}{4}
The governing equation $\vec\tau=\dfrac{d\vec L}{dt}$, for rigid body motion bears formal resemblance to the translational equation $\vec F=\dfrac{d\vec P}{dt}$ where linear momentum $\vec P=M\vec v$ always holds true as $M$ does not have only directional value. On the other hand for rotational about fixed axis writing $\vec L=I\vec\omega$ can not be correct as $I$ has orientational value for rotation of a rigid body about a fixed axis, angular momentum can still be written as $\vec L=I_{xx}\vec\omega_x+I_{yy}\vec\omega_y+I_{zz}\vec\omega_z$ about three mutually perpendicular axes if there axes are symmetrical with respect to body. Symmetry we mean that with chosen coordinates axes $\sum m_ix_iy_i,\ \sum m_iy_iz_i\ \&\ \sum m_ix_iz_i$ are all zero.\par
Such axes are called principal axes. For a body with cylinderical symmetry, axis of revolution is a principal axis and other two principal axes are mutually perpendicular and lie in a plane through centre of mass, perpendicular to axis of rotation. Kinetic energy in pure rotation can then be written as\par\noindent\hfill{\small$K=\dfrac12I_{xx}\omega_x^2+\dfrac12I_{yy}\omega_y^2+\dfrac12I_{zz}\omega_z^2$}\hfill\null\par
The normal to a uniform ring of mass $2.5$ kg and radius $1$ m is making an angle $\theta=45^\circ$ with the vertical axis passing through is centre. The ring is rotating with angular speed $\omega=\sqrt{\dfrac25}$ rad/s about the vertical axis as shown.
\qfig{p2_phy_s3_para3_4_1}
\qn{3} Magnitude of angular momentum of the ring (in $\mathrm{kg\text{-}m^2/s}$)
\qn{4} Magnitude of the torque (in $\mathrm{kg\text{-}m^2/s^2}$) required to keep the ring rotating is
\paraq{5}{6}
A large uniform disk rotates counter clockwise (when viewed from above) with angular speed $\omega_1=1$ rad/s about an axis through its center of mass and perpendicular to the flat surface of the disk. A second uniform disk with the same mass `$m$' as that of first disk and half the radius, rotates clockwise when viewed from above with angular speed $\omega_2=3$ rad/s about an axis through its center of mass and perpendicular to the flat surface of the disk. Suddenly, the second disk is allowed to fall vertically onto the first disk and the two disks finally move with common angular velocity. [$m=10$ kg, $R=1$ m]
\qfig{p2_phy_s3_para5_6_1}
\qn{5} If is the angular velocity of the disks finally after relative motion stops is $\beta$ rad/s then find the value of $49\beta$.
\qn{6} What should be the average value of torque required so that the first disk keeps on rotating with $\omega_1=1$ rad/s during $10$ sec of time of relative motion (after $10$ sec relative motion stops) of the two disks ?
\sect{SECTION-II(ii)}
\qn{7} Consider a thin hollow cylinder of mass $m$, radius $R$ and length $\ell$. $AA'$ is axis through centre and $BB'$ is at a distance $4R$ parallel to $AA'$. Now cylinder is heated uniformly for a long time. Its radius becomes $\dfrac32R$. Change in moment of inertia due to heating along $BB'$ is $\dfrac{nmR^2}{4}$. Find the value of $n$.
\qfig{p2_phy_s4_q07_1}
\qn{8} If a car is moving rightward with acceleration $a=g\sqrt k$ m/s rightward as shown in the figure. Find the value of $k$ so that, rod maintains its orientation as shown in the figure. Neglect the friction and mass of the small rollers at $A$ and $B$.
\qfig{p2_phy_s4_q08_1}
\qn{9} A uniform rod $AB$ of length $4$m and mass $12$ kg is thrown vertically upwards such that just after the projection the centre of mass of the rod moves vertically upwards with a velocity $10$ m/s and at the same time it is rotating with an angular velocity $\pi/2$ rad/sec about a horizontal axis passing through its mid point. Just after the rod is thrown it is horizontal as shown in the figure. Find the acceleration (in $\mathrm{m/s^2}$) of the point $A$ when the centre of mass is at the highest point. (Take $g=10\ \mathrm{m/s^2}$ and $\pi^2=10$)
\qfig{p2_phy_s4_q09_1}
\partsec{2}{CHEMISTRY}
\sect{SECTION-I(i)}
\qn{1} When $1000$ bar constant external pressure is applied on BCC crystal of metal `$M$' then its crystalline structure is changed from BCC to FCC. If radius of metal is $100\sqrt3$ \AA{} and initially in the BCC lattice there were $2\times10^{15}$ unit cells then : [Take $\sqrt6=2.5$]
\begin{qoptions}
\opt{A}{Work done on the metal crystal is $0.3$ J}
\opt{B}{Work done on the metal crystal is $0.7$ J}
\opt{C}{When crystal changes its structure from BCC to FCC its density increases.}
\opt{D}{When crystal changes its structure from FCC to BCC its density remains constant.}
\end{qoptions}
\qn{2} At moderate pressures, the compressibility factor for a particular gas is given by $Z=1+0.2P-\dfrac{60P}{T}$. Where $P$ is in bar and $T$ is in Kelvin. In which of the following conditions the gas behaves ideally :
\begin{qoptions}
\opt{A}{$P=1$ bar, $T=300$ K}
\opt{B}{$P=100$ bar, $T=400$ K}
\opt{C}{$P=2$ bar, $T=300$ K}
\opt{D}{$P=5$ bar, $T=330$ K}
\end{qoptions}
\qn{3} A metal ($M$) is extracted from its sulphide ore ($X$) by using following steps :\\
After sedimentation, ore is agitated by compressed air in the mixture of water and pine oil. Concentrated ore is burnt with excess of oxygen gives resulting product ($p$). ($p$) is either smelted with anthracite to produce metal or ($p$) is treated with dil. $\mathrm{H_2SO_4}$ to dissolve ($p$) and filtered solution is electrolized for the deposition of metal.\\
Which of the following steps is/are involved in the metallurgy of metal ($M$) from metal sulphide($X$).
\begin{qoptions}
\opt{A}{$(X)\to$ Froth Flotation $\to$ Roasting $\to$ Leaching $\to$ Carbon reduction $\to M$}
\opt{B}{$(X)\to$ Froth Flotation $\to$ Roasting $\to$ Carbon reduction $\to M$}
\opt{C}{$(X)\to$ Froth Flotation $\to$ Roasting $\to$ Leaching $\to$ Electrolytic reduction $\to M$}
\opt{D}{$(X)\to$ Sedimentation $\to$ Froth Flotation $\to$ Roasting $\to$ Self reduction $\to M$}
\end{qoptions}
\qn{4} How many of the following statement is/are INCORRECT ?
\begin{qoptions}
\opt{A}{EAN of $\mathrm{[Rh(ox)_3]^{3-}}$ is $56$}
\opt{B}{$d_{\mathrm{C-C}}$ in Zeise's salt is less than that of $\mathrm{C_2F_4}$ molecule}
\opt{C}{Red bauxite consists of $\mathrm{Fe_2O_3}$ as minor impurity}
\opt{D}{In Hall-Heroult process, the produced $\mathrm{Al}$ floats at the top of the electrolyte}
\end{qoptions}
\qn{5} 
\qfig{p2_chem_s1_q05_1}
Select correct statement(s) for above reaction sequence :
\begin{qoptions}
\opt{A}{$A$ is pair of enantiomer and $A$ and $B$ are distereomers}
\opt{B}{$B$ is pair of enantiomer and $A$ and $B$ are distereomers}
\opt{C}{$A$ represents meso compound}
\opt{D}{$B$ represents meso compound}
\end{qoptions}
\qn{6} 
\qfig{p2_chem_s1_q06_1}
which of the following can be formed by above reaction?
\begin{qoptions}
\opt{A}{\qfigi{p2_chem_s1_q06_A}}
\opt{B}{$\mathrm{CH_3-OH}$}
\opt{C}{$\mathrm{CO_2}$}
\opt{D}{\qfigi{p2_chem_s1_q06_D}}
\end{qoptions}
\sect{SECTION-I(ii)}
\paraq{7}{8}
Ellingham diagram provide significant idea for the metallurgical extraction under various temperature range.
\qfig{p2_chem_s2_para7_8_1}
\qn{7} Between $T_1$ and $T_3$ temperatures which of the following is best reducing agent for other metal oxide :-
\begin{qoptions}
\opt{A}{$X$}
\opt{B}{$Y$}
\opt{C}{$Z$}
\opt{D}{Both `$X$' and `$Y$'}
\end{qoptions}
\qn{8} Temperature after which `$Y$' can reduce `$ZO$' is.
\begin{qoptions}
\opt{A}{$T_1$}
\opt{B}{$T_2$}
\opt{C}{$T_3$}
\opt{D}{None of these}
\end{qoptions}
\paraq{9}{10}
For the given scheme of reaction\par
\noindent\hfill{\scriptsize$\mathrm{C_6H_6}\xrightarrow[\mathrm{liq.\,NH_3}]{\mathrm{Na}}\mathrm{(P)}\xrightarrow[\mathrm{(1\ eq.)}]{\mathrm{Ni/H_2}}\mathrm{(Q)}\xrightarrow{\mathrm{Bayer's\ reagent}}\mathrm{(R)}\xrightarrow{\mathrm{HIO_4}}\mathrm{(S)}\xrightarrow{\mathrm{H^+/KMnO_4}}\mathrm{(T)}\xrightarrow{\Delta}\mathrm{(U)}$}\hfill\null\par
\qn{9} Which of the following are incorrect options :
\begin{qoptions}
\opt{A}{$Q$ will give same product with Markovnikov and antimarkovnikov}
\opt{B}{$R$ is vicinal diol}
\opt{C}{$S$ gives positive Tollen's test}
\opt{D}{$U$ can also be obtained by reaction of cyclobutanone with m-CPBA}
\end{qoptions}
\qn{10} Which of the following statements are correct :
\begin{qoptions}
\opt{A}{$2$ moles of $\mathrm{CO_2}$ is released from $T\to U$}
\opt{B}{$3$ moles of $\mathrm{CH_3MgBr}$ are consumed when reacted with $U$}
\opt{C}{$Q$ when reacted with NBS / $h\nu$ total number of monobromo allylic substitution products obtained will be $2$}
\opt{D}{Bayer's reagent is responsible for anti hydroxylation in alkene's}
\end{qoptions}
\sect{SECTION-II(i)}
\paraq{1}{2}
The $\Delta H^\circ_f$ versus temperature plot for phase $\alpha$ and $\beta$ at $1$ bar pressure is given.
\qfig{p2_chem_s3_para1_2_1}
The transition temperature for $\alpha$ to $\beta$ phase change is $600$ K and $C_{p,\beta}-C_{p,\alpha}=1\ \mathrm{J\,mol^{-1}K^{-1}}$. Assume $(C_{p,\beta}-C_{p,\alpha})$ is independent of temperature in the range of $200$ to $700$ K. $C_{p,\alpha}$ and and $C_{p,\beta}$ are molar heat capacities of $\alpha$ and $\beta$ phases, respectively.
\qn{1} The value of enthalpy change, $H_\beta-H_\alpha$ (in $\mathrm{J\,mol^{-1}}$), at $300$ K is ...............
\qn{2} If $20$ moles of $\beta$ phase are converted to $\alpha$ phase at $500$ K and $1$ bar pressure then heat lost to the surroundings (in kJ) is :
\paraq{3}{4}
Tollen's test (Silver mirror test)\\
$\mathrm{RCHO+[Ag(NH_3)_2]OH}$ (Tollen's reagent) $\xrightarrow{\Delta}\mathrm{RCOOH+NH_3+Ag}\downarrow$ (Silver mirror)\\
Aldehyde acts as reducing agent, they can reduce mild oxidizing agents like Tollen's reagent.
\qn{3} Which of the following compounds give positive Tollen's test and fehling test :\\
HCHO, Ph $-$ CHO, HCOOH,\\
\qfigi{p2_chem_s3_q03_1}\\
\qfigi{p2_chem_s3_q03_2}
\qn{4} \qfig{p2_chem_s3_q04_1}
\par\noindent\hfill{\footnotesize$\xrightarrow[\Delta]{t\mathrm{BuO^-}}\mathrm{(A)}\xrightarrow[\mathrm{Zn/H_2O}]{\mathrm{O_3}}\mathrm{(B)}\xrightarrow{[\mathrm{Ag(NH_3)_2}]\mathrm{OH}}\mathrm{(C)}\xrightarrow{\Delta}\mathrm{(D)}\xrightarrow{\mathrm{NaBH_4}}\mathrm{(E)}$}\hfill\null\par
(Major product of A)\\
Number of $-$OH group present in $E$ are :
\paraq{5}{6}
Two moles of an ideal monoatomic gas are taken along the path $ABC$ as shown in following figure. Given $T_B=2T_A$, $T_A=300$ K.
\qfig{p2_chem_s3_para5_6_1}
\qn{5} Increase in internal energy in $AB$ process is (in kJ) :
\qn{6} Magnitude of work done by the gas in process $ABC$ (in kJ) is :
\sect{SECTION-II(ii)}
\qn{7} Find the number of paramagnetic compound(s) with octahedral anion.\\
$\mathrm{K_2[CoF_6]}$;\quad $\mathrm{O_2[AsF_6]}$;\quad $\mathrm{K_2[CoCl_4]}$;\quad $\mathrm{K_3[Co(CN)_6]}$;\\
$\mathrm{K_3[CoF_6]}$;\quad $\mathrm{[Ni(NH_3)_6]SO_4}$;\quad $\mathrm{[Pt(NH_3)_6][Fe(CN)_6]}$;\quad $\mathrm{Na_2[\overset{II}{Fe}(CN)_5NO]}$
\qn{8} Find the number of lone pair present in $\mathrm{C_3N_3(N_3)_3}$.
\qn{9} If hydrolysis of interhalogen compound can be represented by following general reaction:\\
$\mathrm{XY_{n_1}}\xrightarrow{\mathrm{Water}}n_1\mathrm{HY}+\mathrm{HXO_{n_2}}$\\
If given interhalogen compound is polar and non $-$ planar, then calculate value of $n_1+n_2$.
\partsec{3}{MATHEMATICS}
\sect{SECTION-I(i)}
\qn{1} Let the function $p$, $q$ and $r$ defined by $p(x)=x^{\ell nx}$, $q(x)=(\ell nx)^x$, $r(x)=\log_{e^x}x$, then the value of $\dfrac{p'(e^{-1})\cdot r'(e^2)}{q'(e)}$ is (where $y'$ denotes the derivative of $y$)
\begin{qoptions}
\opt{A}{$\dfrac1e$}
\opt{B}{$\dfrac1{e^2}$}
\opt{C}{$\dfrac2e$}
\opt{D}{$\dfrac2{e^2}$}
\end{qoptions}
\qn{2} Let $I=\displaystyle\int\dfrac{dx}{a+b\cos x}$ where $a,b>0$ and $a+b=\lambda_1$ and $a-b=\lambda_2$ then (where $c$ is the constant of integration)
\begin{qoptions}
\opt{A}{If $\lambda_2=0$ then $I=\dfrac{2}{\lambda_1}\tan\dfrac x2+c$}
\opt{B}{If $\lambda_2>0$ then $I=\dfrac{1}{\sqrt{\lambda_1\lambda_2}}\tan^{-1}\left(\sqrt{\dfrac{\lambda_2}{\lambda_1}}\tan\dfrac x2\right)+c$}
\opt{C}{If $\lambda_2<0$ then\par\noindent\hfill{\footnotesize$I=\dfrac{2}{\sqrt{-\lambda_1\lambda_2}}\log_e\left|\dfrac{\sqrt{\lambda_1}+\sqrt{-\lambda_2}\tan\frac x2}{\sqrt{\lambda_1}-\sqrt{-\lambda_2}\tan\frac x2}\right|+c$}\hfill\null\par}
\opt{D}{If $\lambda_2<0$ then\par\noindent\hfill{\footnotesize$I=\dfrac{1}{\sqrt{-\lambda_1\lambda_2}}\log_e\left|\dfrac{\sqrt{\lambda_1}+\sqrt{-\lambda_2}\tan\frac x2}{\sqrt{\lambda_1}-\sqrt{-\lambda_2}\tan\frac x2}\right|+c$}\hfill\null\par}
\end{qoptions}
\qn{3} Let $f(x)$ be a cubic polynomial which is having local maximum at $(1,2)$ and $f'(x)$ has local extremum at $x=0$. If $f(0)=1$ then
\begin{qoptions}
\opt{A}{Number of real roots of $f(x)-f(-x)=6x-10$ is $1$}
\opt{B}{If $g(x)+f(x)=0$ then complete interval in which $g(x)$ is increasing is $x\in(-\infty,-1]\cup[1,\infty)$}
\opt{C}{The range of values of $K$ for which $f(x)=K$ has three distinct solutions is $K\in(0,2)$}
\opt{D}{Range of the function $P(\theta)=f(\sin\theta)$ where $\theta\in R$ is $[0,2]$}
\end{qoptions}
\qn{4} If $f(x)=\displaystyle\lim_{n\to\infty}\dfrac{2x^{2n}+x+2}{x^{2n}-x^2+3}$, then which of the following statement(s) holds good ?
\begin{qoptions}
\opt{A}{$f(x)$ is continuous at $x=1$}
\opt{B}{$f(x)$ is discontinuous at $x=1$}
\opt{C}{$\displaystyle\lim_{x\to2}f(x)$ exists finitely}
\opt{D}{$\displaystyle\lim_{x\to2}f(x)$ does not exist}
\end{qoptions}
\qn{5} Let $a\in R$ and $f:R\to R$ be given by $f(x)=3x^5-50x^3+135x-a$. The equation $f(x)=0$ will have :
\begin{qoptions}
\opt{A}{one real root for $|a|>216$}
\opt{B}{three real roots for $88<|a|<216$}
\opt{C}{five real roots for $|a|<88$}
\opt{D}{five real roots for $a\in R$}
\end{qoptions}
\qn{6} If $\displaystyle\int\dfrac{x^6-x^3-2}{x^3\sqrt{(x^3+1)^2+3x^4}}\,dx=f(x)+C$, then -
\begin{qoptions}
\opt{A}{$f(1)<f(-1)$}
\opt{B}{$f(2)>f(3)$}
\opt{C}{$f(1)>f(-1)$}
\opt{D}{$f(2)<f(3)$}
\end{qoptions}
\sect{SECTION-I(ii)}
\paraq{7}{8}
Let $y=f(x)$ be a fourth degree algebraic polynomial function such that\\
(i) $\displaystyle\lim_{x\to1}\dfrac{f(x)}{(x-1)^2}=2$\\
(ii) It cuts $y$-axis a point where slope of normal is $\dfrac18$\\
(iii) Tangent at $x=2$ is parallel to $y=8x+1$\\
Consider $y=g(x)$ be a quadratic curve touching and cutting $y=f(x)$ on $x$-axis and $y$-axis respectively.
\qn{7} The angle of intersection between $y=f(x)$ \& $y=g(x)$ at point where they cut $y$-axis is -
\begin{qoptions}
\opt{A}{$\tan^{-1}\dfrac{2}{33}$}
\opt{B}{$\tan^{-1}\dfrac{2}{49}$}
\opt{C}{$\tan^{-1}\dfrac{14}{33}$}
\opt{D}{$\tan^{-1}\dfrac{14}{49}$}
\end{qoptions}
\qn{8} The abscissae of point on curve $y=g(x)$, where length of subtangent is $2018$ can be -
\begin{qoptions}
\opt{A}{$1010$}
\opt{B}{$2019$}
\opt{C}{$3028$}
\opt{D}{$4037$}
\end{qoptions}
\paraq{9}{10}
Consider the series\par
\noindent\hfill{\footnotesize$\underbrace{1^2x^1}_{1\text{ term}}+\underbrace{2^2x^2+2^2x^3}_{2\text{ terms}}+$}\hfill\null\par\noindent\hfill{\footnotesize$\underbrace{4^2x^4+4^2x^5+4^2x^6+4^2x^7}_{4\text{ terms}}+$}\hfill\null\par\noindent\hfill{\footnotesize$\underbrace{8^2x^8+8^2x^9+\ldots}_{8\text{ terms}}+\ldots\ldots$}\hfill\null\par
Now answer the following
\qn{9} If $x=1$, then sum of first $1023$ terms is equal to
\begin{qoptions}
\opt{A}{$\dfrac{4^{20}-1}{3}$}
\opt{B}{$\dfrac{8^{10}-1}{7}$}
\opt{C}{$(2^{1023}-1)$}
\opt{D}{$16^5$}
\end{qoptions}
\qn{10} If $x=-1$, then sum of first $1023$ terms is equal to
\begin{qoptions}
\opt{A}{$-1$}
\opt{B}{$\dfrac{2^{20}-1}{3}$}
\opt{C}{$\dfrac{4^{10}-1}{3}$}
\opt{D}{$\dfrac{8^{10}-1}{7}$}
\end{qoptions}
\sect{SECTION-II(i)}
\paraq{1}{2}
Let $f(x)$ be a polynomial function satisfying the following conditions $\displaystyle\lim_{x\to\infty}\dfrac{f(x)}{|x|^3}=0$, $\displaystyle\lim_{x\to\infty}\left(\sqrt{f(x)}-x\right)=-1$ and $f(0)=0$ (Note : $[k]$ denotes greatest integer function less than or equal to $k$)\\
On the basis of above information, answer the following questions :
\qn{1} The number of points of discontinuity of $g(x)=[f(x)]$ in $[0,3]$ is
\qn{2} The number of points of non-derivability of $h(x)=|f(|x|)|$ is -
\paraq{3}{4}
For all $x\in\left[-\dfrac72,\dfrac72\right]$, let $f(x)=\begin{cases}\bigl|[x]\bigr| & ;\ 0\le\{x\}<\dfrac12\\ |x| & ;\ \dfrac12\le\{x\}<1\end{cases}$\\
(where $[k]$ denotes greatest integer function less than or equal to $k$ and $\{k\}$ denotes fractional part of $k$)\\
On the basis of above information, answer the following questions :
\qn{3} Number of solution of the equation $f(x)=\dfrac52$ is equal to
\qn{4} If $L$ denotes number of points of discontinuity and $M$ is the number of points of non-differentiability of function $f(x)$ then $(L+M)$ is equal to
\paraq{5}{6}
Consider a polynomial for $f(x)=a^2x^3+ax^2-x+b$ with $x=-1$ as point of extrema and $f(-1)>0$.\\
On the basis of above information, answer the following questions :
\qn{5} Let $X=\{b:b\in I$ for $a\in I\}$ and $Y=\{b:b\in I$ for $a\notin I\}$. Then the maximum value of\par\noindent\hfill{\small$Z=\displaystyle\sum_{n\in X}\left(\dfrac{\cos\theta}{2}\right)^n+\sum_{m\in Y}\left(\dfrac{\cos\frac\theta2}{3}\right)^m$}\hfill\null\par
\qn{6} The point of intersection of tangents at $(0,f(0))$ and $(f(0),f(f(0)))$ for $f(x)=a^2x^3+ax^2-x+n$ where $a\in I$ and $n$ is the corresponding smallest positive integral value of $b$ is $(x_0,y_0)$, then $(x_0+y_0)$ is
\sect{SECTION-II(ii)}
\qn{7} \par\noindent\hfill{\small$f(x)=\tan\bigl(\sin^{-1}(\cos(\sin^{-1}x))\bigr)$}\hfill\null\par\noindent\hfill{\small$\tan\bigl(\cos^{-1}(\sin(\cos^{-1}x))\bigr)$}\hfill\null\par where $x\in(0,1)$ then\par\noindent\hfill{\small$\displaystyle\sum_{r=0}^{\infty}\dfrac{1}{2^r}f\left(\dfrac{1}{2^r}\right)$}\hfill\null\par is equal to
\qn{8} Let $p_1$ and $p_2$ be two values of $p$ for which $f(x)=\begin{cases}\dfrac{x\left(\ell n(1+x)+\ell n(1-x)\right)}{(\sec x-\cos x)} & ;\ x\in(-1,0)\\[2mm] \left(p^2-3p-1\right)\sin x+x^2 & ;\ x\in[0,\infty)\end{cases}$ is differentiable at $x=0$ find $p_1^2+p_2^2$
\qn{9} A curve is represented by $x=\sec^2t$ and $y=\cot t$, where $t$ is the parameter. If the tangent at point $P$ on curve where $t=\dfrac\pi4$ meets the curve again at $Q$ then $|PQ|=\dfrac{p\sqrt q}{2}$, where $p$, $q$ are coprime. Find $p+q$
\end{multicols}
\end{document}
